% ============================================================
%  part12.tex  |  原文 PDF p.87–92（页脚 p.86–91）
%  内容：第 9 章（续）9.5 整体理论（Global Theory）
%  说明：本片段是第 9 章的续篇，不重复 \section 与 keypoints。
% ============================================================

% >>> p087
\subsection{整体理论（Global Theory）}

回顾局部理论：我们讨论过 TQM 的一阶形式，使得在带有映射
$\varphi:\Omega^\bullet_{S^1}\to V\simeq\mathbb{R}^{2n}$ 的局部模型中，关联映射
\[
  \langle\cdots\rangle_{\mathrm{free}}:\ C_{-\bullet}(\mathcal{W}_{2n})
  \longrightarrow\widehat{\Omega}^{-\bullet}_{2n}((\hbar))
\]
将 $b$ 与 $\hbar\Delta$、$B$ 与 $d_{2n}$ 分别交织。

% [Tutor Note: 记号提醒：$b$ 与 $B$ 是（周期）循环链复形上的两个微分——$b$ 是
% Hochschild 微分、$B$ 是 Connes 算子，满足 $b^2=B^2=0$、$bB+Bb=0$；$C_{-\bullet}$
% 与 $\widehat{\Omega}^{-\bullet}_{2n}$ 的负分次来自 TQM 中形式次数的整体平移
% （准确定义见 9.1–9.4 节）。局部模型中二者的对应关系为
% $b\leftrightarrow\hbar\Delta$、$B\leftrightarrow d_{2n}$，其中 $\hbar$ 因子
% 不可漏掉。]

本节中，我们要把这一构造粘合到一个辛流形上，并把代数指标定理归结为万有 Lie 代数
上同调的计算。基本想法是把局部模型 $\Sigma\to T^{\mathrm{Model}}\subset X$ 粘合
起来 [24]。

\paragraph*{借助 Gelfand--Kazhdan 形式几何的粘合（Gluing via Gelfand--Kazhdan formal geometry）}

\begin{wdefn}{9.14}
  \textbf{Harish--Chandra 对}是一个二元组 $(\mathfrak{g},K)$，其中 $\mathfrak{g}$ 是
  Lie 代数、$K$ 是 Lie 群，并带有
  \begin{itemize}[leftmargin=1.6em,itemsep=0.2em]
    \item $K$ 在 $\mathfrak{g}$ 上的一个作用：$K\xrightarrow{\ \rho\ }\Aut(\mathfrak{g})$；
    \item 一个自然嵌入 $\operatorname{Lie}(K)\xrightarrow{\ i\ }\mathfrak{g}$，其中
          $\operatorname{Lie}(K)$ 是 $K$ 所对应的 Lie 代数，
  \end{itemize}
  使得二者相容：
  \begin{figure}[htbp]
  \centering
  \includegraphics[width=0.30\linewidth]{fig_p087_1.png}
  \caption*{Harish--Chandra 对的相容性交换图（原文 PDF p.87 示意图，原文无编号）}
\end{figure}
\end{wdefn}

% [Tutor Note: 上述相容性的确切意思是 $d\rho=\operatorname{ad}\circ i$，其中
% $\operatorname{ad}:\mathfrak{g}\to\operatorname{Der}(\mathfrak{g})$ 是伴随作用；
% 它是“群作用与 Lie 代数作用相容”的无穷小形式。检验一个 $(\mathfrak{g},K)$ 对
% 是否合法，只需验证这一条交换性。]

\begin{wdefn}{9.15}
  一个 $(\mathfrak{g},K)$-模是一个向量空间 $V$，带有
  \begin{itemize}[leftmargin=1.6em,itemsep=0.2em]
    \item $K$ 在 $V$ 上的一个作用：$K\xrightarrow{\ \varphi\ }\mathrm{GL}(V)$；
    \item 一个 Lie 代数同态 $\mathfrak{g}\to\End(V)$，
  \end{itemize}
  使得二者相容：
  \begin{figure}[htbp]
  \centering
  \includegraphics[width=0.28\linewidth]{fig_p087_2.png}
  \caption*{(g,K)-模的相容性交换图（原文 PDF p.87 示意图，原文无编号）}
\end{figure}
\end{wdefn}

\begin{wdefn}{9.16}
  $X$ 上的一个\textbf{平坦 $(\mathfrak{g},K)$-丛}是指
  \begin{itemize}[leftmargin=1.6em,itemsep=0.2em]
    \item 一个主 $K$-丛 $P\xrightarrow{\ \pi\ }X$；
    \item $P$ 上一个 $K$-等变的 $\mathfrak{g}$-值 1-形式
          $\gamma\in\Omega^1(P,\mathfrak{g})$，
  \end{itemize}
  且满足以下条件：
  \begin{enumerate}[label=(\arabic*),leftmargin=1.8em,itemsep=0.2em]
    \item 对任意 $a\in\operatorname{Lie}(K)$，设 $\xi_a\in\Vect(P)$ 是由 $a$ 生成的
          向量场。则缩并满足 $\gamma(\xi_a)=a$，且下图交换：
      \begin{figure}[htbp]
  \centering
  \includegraphics[width=0.28\linewidth]{fig_p087_3.png}
  \caption*{平坦 (g,K)-丛条件 (1) 的交换图（原文 PDF p.87 示意图，原文无编号）}
\end{figure}
% >>> p088
    \item $\gamma$ 满足 Maurer--Cartan 方程
      \[
        \dd\gamma+\frac{1}{2}[\gamma,\gamma]=0,
      \]
      其中 $\dd$ 是 $P$ 上的 de Rham 微分，$[-,-]$ 是 $\mathfrak{g}$ 中的 Lie 括号。
  \end{enumerate}
\end{wdefn}

% [Tutor Note: 条件 (1) 中 $\gamma(\xi_a)=a$ 表示 $\gamma$ 沿 $K$-轨道方向恰好等于
% 主丛联络的“垂直分量”；条件 (2) 的 Maurer--Cartan 方程是曲率为零的表述。注意
% $\frac12$ 因子与 $[\gamma,\gamma]$ 的约定必须一致：若把 $[\gamma,\gamma]$ 定义为
% $\gamma^a\wedge\gamma^b[T_a,T_b]$，则曲率为 $\dd\gamma+\frac12[\gamma,\gamma]$；
% 有些文献把 $\frac12$ 吸收进括号的定义，核对时务必确认，这是符号错误的高发点。]

给定平坦 $(\mathfrak{g},K)$-丛 $P\to X$ 与 $(\mathfrak{g},K)$-模 $V$，令
\[
  \Omega^\bullet(P,V):=\Omega^\bullet(P)\otimes V
\]
表示 $P$ 上取值于 $V$ 的微分形式。它带有一个联络
\[
  \nabla^\gamma=\dd+\gamma:\ \Omega^\bullet(P,V)\longrightarrow\Omega^{\bullet+1}(P,V),
\]
由 Maurer--Cartan 方程知该联络是平坦的。群 $K$ 作用在 $\Omega^\bullet(P)$ 与 $V$
上，从而在 $\Omega^\bullet(P,V)$ 上诱导一个自然作用。

令
\[
  V_P:=P\times_K V
\]
为 $X$ 上与 $K$-表示 $V$ 相伴的向量丛。设 $\Omega^\bullet(X;V_P)$ 表示 $X$ 上取值
于丛 $V_P\to X$ 的微分形式。与通常的主丛情形类似，$\nabla^\gamma$ 在 $V_P\to X$
上诱导一个平坦联络。这就定义了（de Rham）链复形
$(\Omega^\bullet(X;V_P),\nabla^\gamma)$，其相应的 de Rham 上同调记为
$H^\bullet(X;V_P)$。

下面我们讨论如何把 Lie 代数上同调下降为 $X$ 上的几何对象。

\begin{wdefn}{9.17}
  设 $V$ 是 $(\mathfrak{g},K)$-模。定义 $(\mathfrak{g},K)$ \textbf{相对 Lie 代数
  上链复形} $(C^\bullet_{\mathrm{Lie}}(\mathfrak{g},K;V),\partial_{\mathrm{Lie}})$ 为
  \[
    C^p_{\mathrm{Lie}}(\mathfrak{g},K;V)
    =\operatorname{Hom}_K\left(\wedge^p(\mathfrak{g}/\operatorname{Lie}(K)),V\right).
  \]
  这里 $\operatorname{Hom}_K$ 表示 $K$-等变线性映射。若把
  $C^p_{\mathrm{Lie}}(\mathfrak{g},K;V)$ 视为 Lie 代数上链
  $C^p_{\mathrm{Lie}}(\mathfrak{g};V)$ 的子空间，则 $\partial_{\mathrm{Lie}}$ 就是
  Chevalley--Eilenberg 微分。显式地，对
  $\alpha\in C^p_{\mathrm{Lie}}(\mathfrak{g},K;V)$，
  \[
    \begin{aligned}
      (\partial_{\mathrm{Lie}}\alpha)(a_1\wedge\cdots\wedge a_{p+1})
      ={}&\sum_{i=1}^{p+1}(-1)^{i-1}a_i\cdot
      \alpha\left(a_1\wedge\cdots\wedge\widehat{a_i}\wedge\cdots\wedge a_{p+1}\right)\\
      &+\sum_{i<j}(-1)^{i+j}\alpha\left([a_i,a_j]\wedge\cdots\wedge\widehat{a_i}
      \wedge\cdots\wedge\widehat{a_j}\wedge\cdots\wedge a_{p+1}\right).
    \end{aligned}
  \]
  相应的上同调记为 $H^\bullet_{\mathrm{Lie}}(\mathfrak{g},K;V)$。
\end{wdefn}

% [Tutor Note: Chevalley--Eilenberg 微分的符号约定：第一项 $(-1)^{i-1}$、第二项
% $(-1)^{i+j}$ 分别对应“把 $a_i$ 移到最前面”与“用 $[a_i,a_j]$ 替换两个变量”所需的
% Koszul 符号；不同教材可能相差一个整体符号，但上同调不受影响。另请注意相对复形中
% 除掉的是 $\operatorname{Lie}(K)$ 而非 $\mathfrak{h}$：二者只在
% $K=\mathrm{Sp}_{2n}$、$\operatorname{Lie}(K)=\mathfrak{sp}_{2n}$ 这类情形下才一致，
% 而定义 9.17 的记号 $(\mathfrak{g},K)$ 允许 $K$ 的 Lie 代数严格小于 $\mathfrak{g}$
% 中的另一个子代数 $\mathfrak{h}$。]

给定 $(\mathfrak{g},K)$-模 $V$ 与带有平坦联络 $\gamma\in\Omega^1(P,\mathfrak{g})$
的平坦 $(\mathfrak{g},K)$-丛 $P\to X$。我们可以定义从 $(\mathfrak{g},K)$ 相对 Lie
代数上链复形到 $P$ 上 $V$-值 de Rham 复形的\textbf{下降映射}（descent map）
\[
  \operatorname{desc}:\ \left(C^\bullet_{\mathrm{Lie}}(\mathfrak{g},K;V),
  \partial_{\mathrm{Lie}}\right)\longrightarrow\left(\Omega^\bullet(X;V_P),\nabla^\gamma\right),
  \qquad \alpha\mapsto\alpha(\gamma,\cdots,\gamma),
\]
并诱导上同调的下降映射
\[
  \operatorname{desc}:\ H^\bullet_{\mathrm{Lie}}(\mathfrak{g},K;V)
  \longrightarrow H^\bullet(X;V_P).
\]

% [Tutor Note: 原文此处说 “to $V$-valued de Rham complex on $P$”，但紧随的公式给出
% 的是 $X$ 上的复形 $(\Omega^\bullet(X;V_P),\nabla^\gamma)$；疑为笔误，应读作
% “on $X$”——$V_P$ 是 $X$ 上的丛，$\gamma$ 只用于把 $\alpha$ 沿 $\gamma$ 取值。
% 请核对原书 p.87（PDF p.88）。]

% >>> p089
\paragraph*{Fedosov 联络再探（Fedosov connection revisited）}

回忆（形式）Weyl 代数
\[
  \mathcal{W}_{2n}=\mathbb{R}[[p_i,q^i]]((\hbar)),\qquad
  \mathcal{W}^+_{2n}=\mathbb{R}[[p_i,q^i]][[\hbar]],
\]
带有诱导的 Lie 代数结构，其 Lie 括号定义为
\[
  [f,g]:=\frac{1}{\hbar}[f,g]_\star=\frac{1}{\hbar}(f\star g-g\star f).
\]

% [Tutor Note: 这里指标写作 $p_i$（下标）与 $q^i$（上标），是余切坐标的写法；
% 前文若出现 $q_i$ 只是指标升降的差别，核对公式时以原图为准。Lie 括号定义中的
% $\frac{1}{\hbar}$ 是为了让 $\star$ 的半经典极限恰好给出 Poisson 括号：
% 若采用 $f\star g=fg+\frac{\hbar}{2}\{f,g\}+\cdots$ 的 Moyal 约定，则
% $[f,g]=\{f,g\}+O(\hbar)$，符号取决于 $\star$ 的方向约定，此处极易差一个负号。]

设 $\mathrm{Sp}_{2n}$ 是保持 Poisson 双向量 $\Pi$ 的线性变换构成的辛群。它通过内
自同构作用在 Weyl 代数上。我们可以把 $\mathrm{Sp}_{2n}$ 的 Lie 代数
$\mathfrak{sp}_{2n}$ 与 $\mathbb{R}[p_i,q^i]$ 中的二次多项式等同起来，而
$\mathfrak{sp}_{2n}$ 是 $\mathcal{W}^+_{2n}$ 的 Lie 子代数。$\mathrm{Sp}_{2n}$ 在
$\mathbb{R}^{2n}$ 上的作用诱导出 $\mathrm{Sp}_{2n}$ 在 $\mathcal{W}^+_{2n}$ 上的
作用。因此，$(\mathcal{W}^+_{2n},\mathrm{Sp}_{2n})$ 与
$(\mathcal{W}_{2n},\mathrm{Sp}_{2n})$ 都是 Harish--Chandra 对。

设 $(X,\omega)$ 是辛流形，$F_{\mathrm{Sp}}(X)$ 是辛标架丛。我们有 Weyl 丛
\[
  \mathcal{W}^+_X=F_{\mathrm{Sp}}(X)\times_{\mathrm{Sp}_{2n}}\mathcal{W}^+_{2n},
  \qquad
  \mathcal{W}_X=F_{\mathrm{Sp}}(X)\times_{\mathrm{Sp}_{2n}}\mathcal{W}_{2n}.
\]
考虑 Harish--Chandra 对
\[
  (\bar{\mathfrak{g}},K)=(\mathfrak{g}/Z(\mathfrak{g}),\mathrm{Sp}_{2n}),
\]
其中 $\mathfrak{g}=\mathcal{W}^+_{2n}$，$Z(\mathfrak{g})=\mathbb{R}[[\hbar]]$ 是
$\mathfrak{g}$ 的中心，$Z(\mathfrak{g})\cap\mathfrak{sp}_{2n}=0$。Fedosov 构造了
平坦 $(\bar{\mathfrak{g}},K)$-丛 $F_{\mathrm{Sp}}(X)\to X$，而
$H^0(X;\mathcal{W}^+_X)$ 给出一个形变量子化。取平凡的
$(\bar{\mathfrak{g}},K)$-模 $\mathbb{R}((\hbar))$。则
\[
  \operatorname{desc}:\ C^\bullet_{\mathrm{Lie}}\left(\mathcal{W}^+_{2n},
  \mathfrak{sp}_{2n};\mathbb{R}((\hbar))\right)\longrightarrow\Omega^\bullet_X((\hbar)).
\]
这就是 Gelfand--Fuks 映射。这里
\[
  C^\bullet_{\mathrm{Lie}}\left(\mathcal{W}^+_{2n},\mathfrak{sp}_{2n};
  \mathbb{R}((\hbar))\right)
  \simeq C^\bullet_{\mathrm{Lie}}\left(\mathcal{W}^+_{2n},
  \mathfrak{sp}_{2n}\oplus Z(\mathcal{W}^+_{2n});\mathbb{R}((\hbar))\right).
\]

% [Tutor Note: Gelfand--Fuks（下降）映射的几何含义是：把相对 Lie 代数上链沿
% Fedosov 联络 $\gamma$ 取值，就得到 $X$ 上的微分形式，从而把纯代数的上同调类
% 送到几何对象上。最后一式的等号把 $\mathfrak{sp}_{2n}$ 换成
% $\mathfrak{sp}_{2n}\oplus Z(\mathcal{W}^+_{2n})$，可理解为把中心视为在复形上
% 作用平凡（作者在此采用简化记号）；若严格按
% $\operatorname{Hom}_{\mathfrak{h}}(\wedge^\bullet(\mathfrak{g}/\mathfrak{h}),-)$
% 的字面定义，商空间会变化，阅读时以作者原意为准。]

\paragraph*{示性类（Characteristic classes）}

我们来回顾 Lie 代数上同调中示性类的 Chern--Weil 构造。它们将经 Gelfand--Fuks
映射下降为通常的示性形式。

设 $\mathfrak{g}$ 是 Lie 代数，$\mathfrak{h}\subset\mathfrak{g}$ 是其 Lie 子代数。
设投影映射
\[
  \operatorname{pr}:\ \mathfrak{g}\longrightarrow\mathfrak{h}
\]
是嵌入 $\mathfrak{h}\subset\mathfrak{g}$ 的 $\mathfrak{h}$-等变分裂。一般来说，
$\operatorname{pr}$ 不是从 $\mathfrak{g}$ 到 $\mathfrak{h}$ 的 Lie 代数同态。
$\operatorname{pr}$ 不是 Lie 代数同态这一事实给出
$R\in\operatorname{Hom}(\wedge^2\mathfrak{g},\mathfrak{h})$：
\[
  R(\alpha,\beta)=[\operatorname{pr}(\alpha),\operatorname{pr}(\beta)]_{\mathfrak{h}}
  -\operatorname{pr}[\alpha,\beta]_{\mathfrak{g}},\qquad \alpha,\beta\in\mathfrak{g}.
\]
$\operatorname{pr}$ 的 $\mathfrak{h}$-等变性蕴含
$R\in\operatorname{Hom}_{\mathfrak{h}}(\wedge^2(\mathfrak{g}/\mathfrak{h}),\mathfrak{h})$。
$R$ 称为\textbf{曲率形式}。设 $\operatorname{Sym}^m(\mathfrak{h}^\vee)^{\mathfrak{h}}$
是 $\mathfrak{h}$ 上齐次次数为 $\deg=m$ 的 $\mathfrak{h}$-不变多项式。给定
$P\in\operatorname{Sym}^m(\mathfrak{h}^\vee)^{\mathfrak{h}}$，我们可以通过复合
\[
  P(R):\ \wedge^{2m}\mathfrak{g}\xrightarrow{\ \wedge^m R\ }
  \operatorname{Sym}^m(\mathfrak{h})\xrightarrow{\ P\ }\mathbb{R}
\]
联系一个上链 $P(R)\in C^{2m}_{\mathrm{Lie}}(\mathfrak{g},\mathfrak{h};\mathbb{R})$。

% [Tutor Note: 这里 $R$ 取值于 $\mathfrak{h}$ 而非 $\mathfrak{g}$，且只在
% $\wedge^2(\mathfrak{g}/\mathfrak{h})$ 上有定义；改变分裂 $\operatorname{pr}$
% 会使 $R$ 改变一个 $\mathfrak{h}$-等变的上链，故下文的上同调类与 $\operatorname{pr}$
% 的选取无关——这正是 Chern--Weil 理论中“曲率形式依赖联络、示性类不依赖联络”
% 的 Lie 代数类比。]

% >>> p090
可以验证 $\partial_{\mathrm{Lie}}P(R)=0$，这在上同调
$H^{2m}(\mathfrak{g},\mathfrak{h};\mathbb{R})$ 中定义一个不依赖于
$\operatorname{pr}$ 选取的上同调类 $[P(R)]$。因此我们得到 Chern--Weil 示性映射的
类比
\[
  \chi:\ \operatorname{Sym}^\bullet(\mathfrak{h}^\vee)^{\mathfrak{h}}
  \longrightarrow H^\bullet(\mathfrak{g},\mathfrak{h};\mathbb{R}),
  \qquad P\mapsto\chi(P):=[P(R)].
\]
现在把上述构造应用于
\[
  \mathfrak{g}=\mathcal{W}^+_{2n},\qquad
  \mathfrak{h}=\mathfrak{sp}_{2n}\oplus Z(\mathfrak{g})
\]
的情形。$\mathfrak{g}=\mathcal{W}^+_{2n}$ 中任意元素 $f$ 都可唯一地写成多项式
$f=f(y^i,\hbar)$，其中坐标取为
\[
  (y^1,\cdots,y^n,y^{n+1},\cdots,y^{2n})=(p_1,\cdots,p_n,q^1,\cdots,q^n).
\]
定义 $\mathfrak{h}$-等变投影
\[
  \operatorname{pr}_1(f)=\frac{1}{2}\sum_{i,j}
  \partial_i\partial_j f\Big|_{y=\hbar=0}\,y^iy^j\in\mathfrak{sp}_{2n},
  \qquad
  \operatorname{pr}_3(f)=f|_{y=0}\in Z(\mathfrak{g}).
\]
我们得到相应的曲率
\[
  \begin{aligned}
    R_1&:=[\operatorname{pr}_1(-),\operatorname{pr}_1(-)]
    -\operatorname{pr}_1[-,-]\in\operatorname{Hom}(\wedge^2\mathfrak{g},\mathfrak{sp}_{2n}),\\
    R_3&:=-\operatorname{pr}_3[-,-]\in\operatorname{Hom}(\wedge^2\mathfrak{g},\mathbb{R}[[\hbar]]).
  \end{aligned}
\]

\begin{wrem}{9.18}
  当我们并入向量丛时，可以考虑更一般的情形：此时
  \[
    \mathfrak{g}=\mathcal{W}^+_{2n}+\hbar\left(\mathfrak{gl}(\mathcal{W}^+_{2n})\right),
    \qquad
    \mathfrak{h}=\mathfrak{sp}_{2n}\oplus\hbar\mathfrak{gl}\oplus Z(\mathfrak{g}).
  \]
  其中多出的投影 $\operatorname{pr}_2$ 及其相应的曲率 $R_2$ 分别定义为
  $\hbar\mathfrak{gl}$ 与 $\operatorname{Hom}(\wedge^2,\mathfrak{gl})$ 中的元素。
  值得指出的是，这里所有的 $\operatorname{Hom}$ 都只是 $\mathbb{R}$-线性的映射，
  而不是 $\mathbb{R}[[\hbar]]$-线性的，尽管 $\mathfrak{g}$ 是一个
  $\mathbb{R}[[\hbar]]$-模。
\end{wrem}

% [Tutor Note: 注 9.18 中 “$\operatorname{Hom}(\wedge^2,\mathfrak{gl})$” 省略了
% 第一个变量 $\mathfrak{g}$，应读作 $\operatorname{Hom}(\wedge^2\mathfrak{g},
% \mathfrak{gl})$；$\mathfrak{g}=\mathcal{W}^+_{2n}+\hbar\,\mathfrak{gl}
% (\mathcal{W}^+_{2n})$ 中的加号表示把矩阵值系数并入 Weyl 代数（半直和），并非
% 通常的直和记号的用法。原文照录，特此说明。]

我们现定义 $\widehat{A}$-亏格
\[
  \widehat{A}(\mathfrak{sp}_{2n}):=\left[\det\left(\frac{R_1/2}
  {\sinh(R_1/2)}\right)^{\frac{1}{2}}\right]\in H^\bullet(\mathfrak{g},\mathfrak{h};\mathbb{R}).
\]

\begin{wprop}{9.19}
  在经 Fedosov 联络的下降映射
  $\operatorname{desc}:H^\bullet(\mathfrak{g},\mathfrak{h};\mathbb{R}((\hbar)))
  \to H^\bullet(X)((\hbar))$ 之下，有
  \[
    \operatorname{desc}\left(\widehat{A}(\mathfrak{sp}_{2n})\right)=\widehat{A}(X),
    \qquad
    \operatorname{desc}(R_3)=\omega_\hbar-\hbar\omega.
  \]
\end{wprop}

% [Tutor Note: $\widehat{A}$-亏格中的 $R_1$ 是 $\mathfrak{sp}_{2n}$-值 2-形式，
% $\det$ 与 $\sinh$ 都在矩阵意义下取，外层的 $\frac12$ 次幂与外层方括号（表示
% 上同调类）都是定义的一部分，不可漏写。$\operatorname{desc}(R_3)=\omega_\hbar
% -\hbar\omega$ 中 $\omega_\hbar$ 是 $\hbar$-修正后的辛形式（Fedosov 构造的
% 标准输出，见 [24]），代入指标公式时它决定 $\ee^{-\omega_\hbar/\hbar}$ 的系数。]

\paragraph*{万有迹映射（Universal trace map）}

回忆利用 $\Omega^\bullet_{S^1}\to\mathbb{R}^{2n}$，我们已得到
\[
  \operatorname{Tr}=\int_{BV}\circ\langle-\rangle_{\mathrm{free}}:\
  CC^{\mathrm{per}}_{-\bullet}(\mathcal{W}_{2n})\longrightarrow
  K:=\mathbb{R}((\hbar))[u,u^{-1}].
\]
我们把它写成
\[
  \operatorname{Tr}\in\operatorname{Hom}_K\left(CC^{\mathrm{per}}_{-\bullet}
  (\mathcal{W}_{2n}),K\right).
\]
这是一个 $(\mathcal{W}^+_{2n},\mathrm{Sp}_{2n})$-模。经平坦
$(\mathcal{W}^+_{2n},\mathrm{Sp}_{2n})$-丛 $F_{\mathrm{Sp}}(X)\to X$，我们得到
相伴丛
\[
  E^{\mathrm{per}}:=F_{\mathrm{Sp}}(X)\times_{\mathrm{Sp}_{2n}}
  \operatorname{Hom}_K\left(CC^{\mathrm{per}}_{-\bullet}(\mathcal{W}_{2n}),K\right)
\]
% >>> p091
带有诱导的平坦联络 $\nabla^\gamma$。

回忆 Weyl 丛 $\mathcal{W}(X)=F_{\mathrm{Sp}}(X)\times_{\mathrm{Sp}_{2n}}
\mathcal{W}_{2n}$ 带有平坦联络 $\nabla^\gamma$。我们想把 $\operatorname{Tr}$
粘合到 $X$ 上。记 $\delta$ 为 $\operatorname{Hom}_K\left(CC^{\mathrm{per}}_{-\bullet}
(\mathcal{W}_{2n}),K\right)$ 上由 $b+uB$ 诱导的微分。于是
\[
  \delta\operatorname{Tr}=\operatorname{Tr}\left((b+uB)(-)\right)=0.
\]
我们可以把 $\operatorname{Tr}$ 看作定义了
\[
  C^0_{\mathrm{Lie}}\left(\mathfrak{g},\mathfrak{h};\operatorname{Hom}_K
  \left(CC^{\mathrm{per}}_{-\bullet}(\mathcal{W}_{2n}),K\right)\right)
\]
中的一个元素，其中取
\[
  \mathfrak{g}=\mathcal{W}^+_{2n}/Z(\mathcal{W}^+_{2n}),\qquad
  \mathfrak{h}=\mathfrak{sp}_{2n}.
\]
然而 $\operatorname{Tr}$ 不是 $\mathfrak{g}$-不变的，即
$\partial_{\mathrm{Lie}}\operatorname{Tr}\neq0$。换言之，$\operatorname{Tr}$ 不是
$(\mathfrak{g},\mathrm{Sp}_{2n})$-模之间的映射。因此 $\operatorname{Tr}$ 不能
直接粘合。

% [Tutor Note: $\delta$ 是由 $b+uB$ 在 $\operatorname{Hom}_K(CC^{\mathrm{per}}_{-\bullet}
% (\mathcal{W}_{2n}),K)$ 上诱导的微分；$b+uB$ 正是周期循环链复形的微分，故
% $\delta\operatorname{Tr}=0$ 表示 $\operatorname{Tr}$ 是（取值于 $K$ 的）循环链的
% 迹。原文 “$\partial_{\mathrm{Lie}}\operatorname{Tr}=\delta(-)$” 是简写，意为
% $\partial_{\mathrm{Lie}}\operatorname{Tr}$ 是某个元素的 $\delta$-像，而不是等于
% $\delta$ 算子本身；照录原文，阅读时注意。]

我们注意到 $\partial_{\mathrm{Lie}}\operatorname{Tr}=\delta(-)$。事实表明，存在
一种典范的方式把 $\operatorname{Tr}$ 提升为
\[
  \widehat{\operatorname{Tr}}\in C^\bullet_{\mathrm{Lie}}\left(\mathfrak{g},\mathfrak{h};
  \operatorname{Hom}_K\left(CC^{\mathrm{per}}_{-\bullet}(\mathcal{W}_{2n}),K\right)\right),
\]
使得
\[
  \widehat{\operatorname{Tr}}=\operatorname{Tr}
  +\ C^{>0}_{\mathrm{Lie}}\left(\mathfrak{g},\mathfrak{h};\operatorname{Hom}_K
  \left(CC^{\mathrm{per}}_{-\bullet}(\mathcal{W}_{2n}),K\right)\right)\ \text{中的项},
\]
并满足耦合上闭链条件
\[
  (\partial_{\mathrm{Lie}}+\delta)\widehat{\operatorname{Tr}}=0.
\]
$\widehat{\operatorname{Tr}}$ 称为\textbf{万有迹映射}。让我们代入
$1\in\mathcal{W}_{2n}$，则 $\widehat{\operatorname{Tr}}(1)$ 是
$\partial_{\mathrm{Lie}}$-闭的，从而定义\textbf{万有指标}
\[
  \left[\widehat{\operatorname{Tr}}(1)\right]
  \in H^\bullet_{\mathrm{Lie}}(\mathfrak{g},\mathfrak{h};K).
\]

\begin{wthm}{9.20}（万有代数指标定理（Universal algebraic index theorem））
  \[
    \left[\widehat{\operatorname{Tr}}(1)\right]
    =u^ne^{-R_3/(u\hbar)}\widehat{A}(\mathfrak{sp}_{2n})_u,
  \]
  其中对 $A=\sum_{p\ \text{even}}A_p$，$A_p\in H^p(\mathfrak{g},\mathfrak{h};K)$，
  定义
  \[
    A_u=\sum_p u^{-p/2}A_p.
  \]
\end{wthm}

% [Tutor Note: 定理中的 $u^n$ 与 $A_u=\sum_p u^{-p/2}A_p$ 表明该指标是 $u$ 的半整数
% 次幂的 Laurent 级数；$e^{-R_3/(u\hbar)}$ 的分母 $u\hbar$ 不可与下文
% $e^{-\omega_\hbar/\hbar}$ 的分母混淆。此处 $[\cdot]$ 表示相对 Lie 代数上同调类；
% 整体化（PDF p.92）后即得 $\int_X\ee^{-\omega_\hbar/\hbar}\widehat{A}(X)$，这正是
% Atiyah--Singer 指标定理的 $\widehat{A}$-形式在形变量子化中的类比。]

该定理由 Feigin--Tsygan [20]、Feigin--Felder--Shoikhet [19]、
Bressler--Nest--Tsygan [8] 等人的工作发展而来，并可自然地推广到向量丛的情形 [24]。

现在对 $\widehat{\operatorname{Tr}}$ 应用 Gelfand--Fuks（下降）映射，使得
\begin{figure}[htbp]
  \centering
  \includegraphics[width=0.45\linewidth]{fig_p091_1.png}
  \caption*{Gelfand--Fuks 下降映射的纵向箭头（原文 PDF p.91 示意图，原文无编号）}
\end{figure}
设 $\mathcal{W}_D(X)$ 是 $\mathcal{W}(X)$ 的平坦截面空间，它给出一个形变量子化。
则
\[
  \operatorname{desc}(\widehat{\operatorname{Tr}}):\
  CC^{\mathrm{per}}_{-\bullet}(\mathcal{W}_D(X))\longrightarrow
  \Omega^\bullet(X)((\hbar))[u,u^{-1}],\qquad b+uB\mapsto \dd_X.
\]

% >>> p092
特别地，它通过
\[
  f\in\mathcal{W}_D(X)\longmapsto\int_X\operatorname{desc}
  (\widehat{\operatorname{Tr}})(f)\in\mathbb{R}((\hbar))
\]
在形变量子化中定义一个迹映射。可以证明
$\int_X\operatorname{desc}(\widehat{\operatorname{Tr}})(f)$ 不涉及 $u$。由万有
代数指标定理，我们有
\[
  \int_X\operatorname{desc}(\widehat{\operatorname{Tr}})(1)
  =\int_X\ee^{-\omega_\hbar/\hbar}\widehat{A}(X).
\]
这就给出\textbf{代数指标定理}。

% [Tutor Note: “$\int_X\operatorname{desc}(\widehat{\operatorname{Tr}})(f)$ 不涉及
% $u$” 是说积分后 $u$ 的依赖消失（相应的项为恰当形式/整体上消失），故得到与 $u$
% 无关的数值不变量；这是把循环同调的迹下降为普通数值迹的关键一步。整体公式
% $\int_X\ee^{-\omega_\hbar/\hbar}\widehat{A}(X)$ 与代数指标定理的标准形式一致，
% $\hbar\to0$ 时给出 $\int_X\widehat{A}(X)$ 这一拓扑指标。]

\paragraph*{万有迹映射 $\widehat{\operatorname{Tr}}$ 的构造（Construction of universal trace map $\widehat{\operatorname{Tr}}$）}

我们有如下关系。
\begin{figure}[htbp]
  \centering
  \includegraphics[width=0.75\linewidth]{fig_p092_1.png}
  \caption*{背景对称性、联络形式与相互作用之间的关系（原文 PDF p.92 示意图，原文无编号）}
\end{figure}
设 $\Theta:\mathfrak{g}\to\mathcal{W}^+_{2n}/Z(\mathcal{W}^+_{2n})
=\bar{\mathfrak{g}}$ 是典范恒等映射。对每个
$f\in\mathcal{W}^+_{2n}/Z(\mathcal{W}^+_{2n})$，我们已在
$\mathcal{E}=\Omega^\bullet(S^1)\otimes\mathbb{R}^{2n}$ 上定义局部泛函
\[
  I_f(\varphi)=\int_{S^1}f(\varphi),\qquad \varphi\in\mathcal{E}.
\]
于是 $\Theta$ 给出映射
\[
  I_\Theta:\ \mathfrak{g}\longrightarrow\mathcal{O}_{\mathrm{loc}}(\mathcal{E}),
  \qquad f\mapsto I_{\Theta(f)}.
\]
我们可以把这个映射看作
\[
  I_\Theta\in C^1(\mathfrak{g},\mathcal{O}_{\mathrm{loc}}(\mathcal{E}))
  =\mathfrak{g}^\vee\otimes\mathcal{O}_{\mathrm{loc}}(\mathcal{E}).
\]
现在可以构造
$\widehat{\operatorname{Tr}}\in C^\bullet_{\mathrm{Lie}}\left(\mathfrak{g},
\mathfrak{h};\operatorname{Hom}_K\left(CC^{\mathrm{per}}_{-\bullet}
(\mathcal{W}_{2n}),K\right)\right)$ 如下：
\[
  \begin{aligned}
    \widehat{\operatorname{Tr}}(f_0\otimes f_1\otimes\cdots\otimes f_m)
    &:=\int_{BV}\exp\left(\hbar P_0^\infty\right)\left(\mathcal{O}_{f_0,f_1,\cdots,f_m}
    \ee^{\frac{1}{\hbar}I_\Theta}\right)
    \in C^\bullet(\mathfrak{g},\mathfrak{h};K),\qquad f_i\in\mathcal{W}_{2n}\\
    &\text{``}=\int_{BV}\int_{\operatorname{Im}\dd^*\subset\mathcal{E}}
    \ee^{-\frac{1}{2\hbar}\int_{S^1}\langle\varphi,\dd\varphi\rangle
    +\frac{1}{\hbar}I_\Theta}\mathcal{O}_{f_0,f_1,\cdots,f_m}\text{''}.
  \end{aligned}
\]

% [Tutor Note: 原文 “Let $\Theta:\mathfrak{g}\to\mathcal{W}^+_{2n}/Z(\mathcal{W}^+_{2n})
% =\bar{\mathfrak{g}}$ be the canonical identity map” 中，$\Theta$ 严格说不是恒等映射
% （定义域与值域不同），而是把 $\mathfrak{g}=\mathcal{W}^+_{2n}$ 等同于商
% $\bar{\mathfrak{g}}$ 的典范等同/商映射；照录原文记号，阅读时按“商映射”理解。
% 第二行首尾的引号表示这是经过 Gauss 积分（对 $\operatorname{Im}d^*\subset
% \mathcal{E}$ 上的场）之后的形式表达式，并非逐字恒等；$P_0^\infty$ 是前文
% （9.1–9.4 节）同伦 RG 流中的传播子。]

\paragraph*{指标的计算（Computation of index）}

Weyl 代数 $\mathcal{W}_{2n}$ 可以看作以 $\hbar$ 为参数的一族结合代数。这导致
\textbf{Gauss--Manin--Getzler 联络} $\nabla_{\hbar\partial_\hbar}$ 作用在
$CC^{\mathrm{per}}_{-\bullet}(\mathcal{W}_{2n})$ 上。指标的计算由以下步骤组成：
\begin{enumerate}[label=(\arabic*),leftmargin=1.8em,itemsep=0.3em]
  \item \emph{Feynman 图计算}给出
    \[
      \widehat{\operatorname{Tr}}(1)=u^ne^{-R_3/(u\hbar)}
      \left(\underbrace{\widehat{A}(\mathfrak{sp}_{2n})_u}_{\text{单圈计算}}
      +O(\hbar)\right).
    \]
  \item \emph{Gauss--Manin--Getzler 联络}的计算表明
    $\nabla_{\hbar\partial_\hbar}\left(\ee^{R_3/(u\hbar)}
    \widehat{\operatorname{Tr}}(1)\right)$ 是 $\partial_{\mathrm{Lie}}$-恰当的。
  \item 结合 (1) 与 (2)，我们得到
    \[
      \left[\widehat{\operatorname{Tr}}(1)\right]
      =\left[u^ne^{-R_3/(u\hbar)}\widehat{A}(\mathfrak{sp}_{2n})_u\right]
      \in H^\bullet(\mathfrak{g},\mathfrak{h};K).
    \]
\end{enumerate}

% [Tutor Note: 第 (1) 步下括号“单圈计算”指出：$\widehat{A}$-亏格来自单圈 Feynman
% 图的和，而高圈项 $O(\hbar)$ 不改变上同调类；第 (2) 步用 Gauss--Manin--Getzler
% 联络（Getzler 在热核证明中引入的联络在循环同调上的版本）说明
% $\ee^{R_3/(u\hbar)}\widehat{\operatorname{Tr}}(1)$ 的协变导数恰当，从而第 (1) 步
% 的结果可以代表整个上同调类。三步合起来正是“局部单圈计算 + 整体下降 = 指标定理”
% 的标准论证结构。]
